\clearpage
\section{Colored complementary
reserve}\label{h-colored-complementary-reserve}

Technical derivation for review. Equation and subsection numbers in this
appendix are local to the source derivation. Historical local labels are
retained, including numbering gaps or nonsequential labels. The
accompanying source notes retain the original expressions for
comparison.

\subsection{1. Coordinate conventions and exact
matrices}\label{h-coordinate-conventions-and-exact-matrices}

Use the intrinsic geometry note and the notation of the actual adapted
source inventory:

% DISPLAY h.1
% SOURCE F=E*B_(Z+A), C=P_NB_A, L=E*B_Y, N_Y=P_NB_Y,
% SOURCE I_Y=P_IB_Y, C_Y=P_CB_Y, K=R^−1C_Y*, J2=KC_Y,
% SOURCE M=F+L−J2, W=I+KK*, g_C=I+K*K.
\begin{gather*}
\begin{aligned}
& F=E^{*}B_{Z+A}, C=P_{N} B_{A}, L=E^{*}B_{Y}, N_{Y}=P_{N} B_{Y}, \\
& I_{Y}=P_{I} B_{Y}, C_{Y}=P_{C} B_{Y}, K=R^{-1}C_{Y}^{*}, J_{2}=KC_{Y}, \\
& M=F+L-J_{2}, W=I+KK^{*}, g_{C}=I+K^{*}K.
\end{aligned}\notag
\end{gather*}

Let \(p_{s}=(g_{C}^{-1/2} p_{C},p_{I})\), where \(p_{C}\) includes the
exact normal connection. Define

% DISPLAY h.2
% SOURCE Ψ=[(C_Y*+(F+L)*K)g_C^−1/2,I_Y*],
% SOURCE A=W^1/2M^−*Ψ,
% SOURCE a=Bp_F=W^1/2M^−*(C+N_Y)*p_F,
% SOURCE s=−W^1/2M^−*S_f.
\begin{gather*}
\begin{aligned}
& \Psi  =[(C_{Y}^{*}+(F+L)^{*}K)g_{C}^{-1/2},I_{Y}^{*}], \\
& A=W^{1/2}M^{-*} \Psi  , \\
& a=Bp_{F}=W^{1/2}M^{-*}(C+N_{Y})^{*}p_{F}, \\
& s=-W^{1/2}M^{-*}S_{f}.
\end{aligned}\notag
\end{gather*}

The minus sign in \(s\) matches the original Gauss expression
\(S_{f}-\mathrm{fast}-\mathrm{slow}\) after multiplying that whole stack
by \(-1\). The unshifted scalar kinetic stack is exactly

% DISPLAY h.3
% SOURCE |p_s v|²+|p_Fv|²+|a v+A p_s v+s v|².                 (1.1)
\begin{gather*}
\begin{aligned}
& \vert p_{s} v\vert ^{2}+\vert p_{F} v\vert ^{2}+\vert a v+A p_{s} v+s v\vert ^{2}.
\end{aligned}\tag{1.1}
\end{gather*}

The scalar half-density term is added after this first-order identity.
We never add the same Gauss square a second time to
\(H_{\mathrm{fast}}\).

Throughout, \(T_{\mathrm{cent}}[v]=\Vert p_{C}v\Vert ^{2}\) with the
exact normal connection. Its direct induced-metric coefficient obeys
\(g_{C}^{-1}\ge (1+C \sigma ^{2})^{-1}I\); this small center-only loss
is included in the final allocation.

Set
\(\mathrm{D0}=F^{-1}C^{*}p_{F}, P=L-J_{2}, Z=F^{-1}P^{*}, T=(I+Z)^{-1}, n=F^{-1}N_{Y}^{*}p_{F}\).
Thus

% DISPLAY h.4
% SOURCE a=W^1/2 T(D0+n),   ||T||+||T^−1||≤C,   ||A||≤Cσ.   (1.2)
\begin{gather*}
\begin{aligned}
& a=W^{1/2} T(\mathrm{D0}+n), \quad  \Vert T\Vert +\Vert T^{-1}\Vert \le C, \quad  \Vert A\Vert \le C \sigma  .
\end{aligned}\tag{1.2}
\end{gather*}

The frozen kinetic contribution in \(H_{\mathrm{fast}}\) is
\(p_{F}^{2}+\Vert \mathrm{D0}\Vert ^{2}\).

\subsection{2. Oscillator estimates with the vacuum kept
explicitly}\label{h-oscillator-estimates-with-the-vacuum-kept-explicitly}

The adapted bosonic Gaussian on edge \(e\) is
\(\exp (-Y_{e} \cdot  \Omega _{e}Y_{e}/2)\), with
\(c r_{e} I\le  \Omega _{e}\le C r_{e}\) I. Let
\(a_{e}= \partial  _{e}+ \Omega _{e}Y_{e}\) be its annihilation
gradient. Then \(\Vert a_{e} v\Vert ^{2}\le C h_{e}[v]\), including
arbitrary fermionic and spectator components. The shifted fermionic
excitation form is nonnegative and only improves this inequality.

The following estimates are consequences of splitting each edge into its
adapted full vacuum \(P_{e}\) and \(Q_{e}=1-P_{e}\); see Appendix \(E\)
for the first two:

% DISPLAY h.5
% SOURCE ||χY_ep_e/r_e v||²≤Cσ²h_e[v]+Cr_e^−2||v||²,         (2.1)
% SOURCE 
% SOURCE ||χY_fp_e/r_g v||²≤C(σ²+r_*^−3)(h_e+h_f)[v]
% SOURCE                          +CΣ_(u in triangle)r_u^−2||v||². (2.2)
\begin{gather*}
\begin{aligned}
& \Vert  \chi  Y_{e} p_{e}/r_{e} v\Vert ^{2}\le C \sigma ^{2}h_{e}[v]+Cr_{e}^{-2}\Vert v\Vert ^{2},
\end{aligned}\tag{2.1} \\
\begin{aligned}
& \Vert  \chi  Y_{f} p_{e}/r_{g} v\Vert ^{2}\le C( \sigma ^{2}+r_{*}^{-3})(h_{e}+h_{f})[v] \\
& +C \sum _{u \in  \text{ triangle }}r_{u}^{-2}\Vert v\Vert ^{2}.
\end{aligned}\tag{2.2}
\end{gather*}

A useful sharpened two-coordinate version is

% DISPLAY h.6
% SOURCE ||χY_fY_gv||²≤Cσ²r_*²(h_f[v]/r_f²+h_g[v]/r_g²)
% SOURCE                          +C/(r_fr_g)||v||²,               (2.3)
\begin{gather*}
\begin{aligned}
& \Vert  \chi  Y_{f} Y_{g} v\Vert ^{2}\le C \sigma ^{2}r_{*}^{2}(h_{f}[v]/r_{f}^{2}+h_{g}[v]/r_{g}^{2}) \\
& +C/(r_{f} r_{g})\Vert v\Vert ^{2},
\end{aligned}\tag{2.3}
\end{gather*}

for distinct edges f,g. To prove it, decompose into
\(Q_{f}, P_{f}Q_{g}\) and \(P_{f}P_{g}\). On the first two pieces use
\(\vert  \chi  Y_{g}\vert \le C \sigma  r_{*}\) or
\(\vert  \chi  Y_{f}\vert \le C \sigma  r_{*}\), and
\(\Vert Y_{f}Q_{f} v\Vert ^{2}\le C h_{f}[v]/r_{f}^{2}\) (similarly g).
On the last use the exact Gaussian second moments. Finite-sum
Cauchy--Schwarz supplies the displayed constant. The cutoff acts only on
the left; it need not commute with \(P_{f}\).

In every triangle e,f,g,

% DISPLAY h.7
% SOURCE r_*r_g/(r_er_f)≤2,    r_g/(r_e²r_f)≤C(r_e^−2+r_f^−2). (2.4)
\begin{gather*}
\begin{aligned}
& r_{*}r_{g}/(r_{e} r_{f})\le 2, \quad  r_{g}/(r_{e}^{2}r_{f})\le C(r_{e}^{-2}+r_{f}^{-2}).
\end{aligned}\tag{2.4}
\end{gather*}

These elementary triangle inequalities are what remove remote zero-point
energies.

\subsection{3. Direct cubic potential, with arbitrary fast
vectors}\label{h-direct-cubic-potential-with-arbitrary-fast-vectors}

Let e,f,g be a triangle, \(s_{t}=r_{e}+r_{f}+r_{g}\), and let \(B_{t}\)
be any fixed bounded trilinear contraction of the three real fast
coordinate spaces, also allowing a bounded spectator matrix. Then

% DISPLAY h.8
% SOURCE |<v,s_t B_t(Y_e,Y_f,Y_g)v>|
% SOURCE   ≤Cσ(h_e+h_f+h_g)[v]+Cσ^−1Σ_(u in t)r_u^−2||v||². (3.1)
\begin{gather*}
\begin{aligned}
& \vert \langle v,s_{t} B_{t}(Y_{e},Y_{f},Y_{g})v\rangle \vert \\
& \le C \sigma  (h_{e}+h_{f}+h_{g})[v]+C \sigma ^{-1} \sum _{u \in  t}r_{u}^{-2}\Vert v\Vert ^{2}.
\end{aligned}\tag{3.1}
\end{gather*}

Choose \(e\) with largest \(r_{e}\). For a scalar component, integration
by parts gives exactly

% DISPLAY h.9
% SOURCE < v,Y_eY_fY_g v >
% SOURCE  =Re< a_ev, Ω_e^−1(Y_fY_g)v >.                           (3.2)
\begin{gather*}
\begin{aligned}
& \langle v,Y_{e} Y_{f} Y_{g} v\rangle \\
& =\operatorname{Re} \langle a_{e} v, \Omega _{e}^{-1}(Y_{f} Y_{g})v\rangle .
\end{aligned}\tag{3.2}
\end{gather*}

There is no derivative of \(Y_{f} Y_{g}\) in the \(e\) coordinate.
Equation (3.2) applies componentwise to the finite contraction, with
\(\Omega _{e}^{-1}\) acting on its \(e\) index. Use
\(s_{t}/r_{e}\le 3\), (2.3), \(r_{*}\le r_{f},r_{g}\) and Young with
parameter \(\sigma\). The source term is at most
\(C \sigma ^{-1}/(r_{f} r_{g})\), which is bounded by the triangle
\(W_{2}\). This proves (3.1). No finite oscillator truncation is made.

Every actual cubic potential term has a bounded coefficient times
\(s_{t}\): its coefficient is a commutator with \(Z+A\), whose norm is
\(\le C(r_{e}+r_{f}+r_{g})\) under the pairwise-relative internal tube.
Therefore (3.1) applies to the full \(V_{3}\). \(V_{4}\) remains
positive and is not spent.

\subsection{4. Signed triangle kinetic
forms}\label{h-signed-triangle-kinetic-forms}

For any bounded real coefficient tensor independent of the fast
variables \(Y\) (but possibly dependent on the slow variables), and
three distinct triangle edges,

% DISPLAY h.10
% SOURCE |2Re<p_ev, (Y_f/r_e)p_gv>|
% SOURCE  ≤Cσ(h_e+h_f+h_g)[v]+Cσ^−1Σ_(u in t)r_u^−2||v||². (4.1)
\begin{gather*}
\begin{aligned}
& \vert 2\operatorname{Re} \langle p_{e} v, (Y_{f}/r_{e})p_{g} v\rangle \vert \\
& \le C \sigma  (h_{e}+h_{f}+h_{g})[v]+C \sigma ^{-1} \sum _{u \in  t}r_{u}^{-2}\Vert v\Vert ^{2}.
\end{aligned}\tag{4.1}
\end{gather*}

The clean proof is the Gaussian ground-state transform. The real
symmetric principal coefficient \(a_{eg}(Y)=Y_{f}/r_{e}\) has no
divergence in either active derivative coordinate e,g. For the product
Gaussian \(g_{B}\),

% DISPLAY h.11
% SOURCE ∫(∂v)*a∂v
% SOURCE  =∫g_B²(∂(v/g_B))*a∂(v/g_B)
% SOURCE   +∫[div(aΩY)−(ΩY)·a(ΩY)]|v|².                         (4.2)
\begin{gather*}
\begin{aligned}
& \int  ( \partial  v)^{*}a \partial  v \\
& = \int  g_{B}^{2}( \partial  (v/g_{B}))^{*}a \partial  (v/g_{B}) \\
& + \int  [\operatorname{div} (a \Omega  Y)-( \Omega  Y) \cdot  a( \Omega  Y)]\vert v\vert ^{2}.
\end{aligned}\tag{4.2}
\end{gather*}

The first term is \(\le C \sigma  \Sigma  h_{e}\) because
\(\vert Y_{f}\vert /r_{e}\le C \sigma\) on the support. The divergence
term is zero since the three roots are distinct. The last term is a
triangle cubic with coefficient bounded by \(C r_{g}\), hence by
\(C s_{t}\), and (3.1) proves (4.1). An outer cutoff equal to one near
the support contributes no derivative term there.

One may instead expand \(p=-i a+i \Omega  Y\). Formula (2.3), together
with (2.4), controls the apparently dangerous factor \(r_{g}/r_{e}\) in
the mixed annihilator/quadratic term. This gives the same bound without
ever estimating \(p_{e}^{2}\) by \(h_{e}\) plus an unassigned \(r_{e}\)
vacuum mass.

\subsection{5. Exact inverse-M reduction: higher paths are products of
errors}\label{h-exact-inverse-m-reduction-higher-paths-are-products-of-errors}

Let \(X=n-Z\mathrm{D0}\). The connection bounds imply

% DISPLAY h.12
% SOURCE ||Xv||²+||Z*D0v||²
% SOURCE  ≤C(σ²+r_*^−3)H_exc[v]+CW2[v].                          (5.1)
\begin{gather*}
\begin{aligned}
& \Vert Xv\Vert ^{2}+\Vert Z^{*}\mathrm{D0}v\Vert ^{2} \\
& \le C( \sigma ^{2}+r_{*}^{-3})H_{\mathrm{exc}}[v]+CW_{2}[v].
\end{aligned}\tag{5.1}
\end{gather*}

For \(X\) this is the displayed helper lemma. For
\(Z^{*}\mathrm{D0}=P F^{-1}\mathrm{D0}\) the linear \(L\) term is a
triangle map \(Y_{f} p_{e}/r_{e}\); the \(J_{2}\) term factors through
\(C_{Y} F^{-1}\mathrm{D0}\), a same-edge map, followed by bounded \(K\).
Thus the identical proof applies. Denominators stay on their actual
input or output root.

Since \(W\ge I\) and \(T-I=-ZT\) exactly,

% DISPLAY h.13
% SOURCE ||av||²≥||D0v+TXv||²
% SOURCE  ≥||D0v||²+2Re<D0v,Xv>−2|<Z*D0v,TXv>|.                 (5.2)
\begin{gather*}
\begin{aligned}
& \Vert av\Vert ^{2}\ge \Vert \mathrm{D0}v+TXv\Vert ^{2} \\
& \ge \Vert \mathrm{D0}v\Vert ^{2}+2\operatorname{Re} \langle \mathrm{D0}v,Xv\rangle -2\vert \langle Z^{*}\mathrm{D0}v,TXv\rangle \vert .
\end{aligned}\tag{5.2}
\end{gather*}

The last term is bounded by (5.1). The \(J_{2}\) part of the remaining
cross factors as

% DISPLAY h.14
% SOURCE <D0,F^−1J2*D0>=<C_YF^−1D0,K*D0>,                       (5.3)
\begin{gather*}
\begin{aligned}
& \langle \mathrm{D0},F^{-1}J_{2}^{*}\mathrm{D0}\rangle =\langle C_{Y} F^{-1}\mathrm{D0},K^{*}\mathrm{D0}\rangle ,
\end{aligned}\tag{5.3}
\end{gather*}

and both factors satisfy the same-edge bound. The two other terms are
precisely

% DISPLAY h.15
% SOURCE 2Re<D0,F^−1N_Y*p_F>−2Re<D0,F^−1L*D0>.                   (5.4)
\begin{gather*}
\begin{aligned}
& 2\operatorname{Re} \langle \mathrm{D0},F^{-1}N_{Y}^{*}p_{F}\rangle -2\operatorname{Re} \langle \mathrm{D0},F^{-1}L^{*}\mathrm{D0}\rangle .
\end{aligned}\tag{5.4}
\end{gather*}

Their root supports are triangles; their coefficients are bounded by
\(C \kappa  /r_{e}\) and \(C \kappa ^{2}/r_{e}\). They obey (4.1).
Consequently

% DISPLAY h.16
% SOURCE ||av||²≥||D0v||²−C(κσ+σ²+r_*^−3)H_exc[v]
% SOURCE                          −C(1+κ/σ)W2[v].                 (5.5)
\begin{gather*}
\begin{aligned}
& \Vert av\Vert ^{2}\ge \Vert \mathrm{D0}v\Vert ^{2}-C( \kappa  \sigma  + \sigma ^{2}+r_{*}^{-3})H_{\mathrm{exc}}[v] \\
& -C(1+ \kappa  / \sigma  )W_{2}[v].
\end{aligned}\tag{5.5}
\end{gather*}

This is an actual all-vector signed-form comparison for the coupled
inverse-M fast column. It does not require a separate estimate of every
higher Neumann path. Importantly, it never bounds the large
\(\mathrm{D0}\) vacuum norm by the nearest complementary gap.

\subsection{6. Exact slow completion and finite-spin
payments}\label{h-exact-slow-completion-and-finite-spin-payments}

With \(s\) temporarily omitted, the identity

% DISPLAY h.17
% SOURCE ||p_s v||²+||a v+A p_s v||²
% SOURCE  =||(p_s+A*a)v||²+||av||²+||A p_sv||²−||A*a v||²         (6.1)
\begin{gather*}
\begin{aligned}
& \Vert p_{s} v\Vert ^{2}+\Vert a v+A p_{s} v\Vert ^{2} \\
& =\Vert (p_{s}+A^{*}a)v\Vert ^{2}+\Vert av\Vert ^{2}+\Vert A p_{s} v\Vert ^{2}-\Vert A^{*}a v\Vert ^{2}
\end{aligned}\tag{6.1}
\end{gather*}

holds pointwise on the first-order stack. The final negative term is
controlled by

% DISPLAY h.18
% SOURCE ||A*a v||²≤C(σ²+r_*^−3)H_exc[v]+CW2[v],                  (6.2)
\begin{gather*}
\begin{aligned}
& \Vert A^{*}a v\Vert ^{2}\le C( \sigma ^{2}+r_{*}^{-3})H_{\mathrm{exc}}[v]+CW_{2}[v],
\end{aligned}\tag{6.2}
\end{gather*}

because
\(A^{*}a=A^{*}\mathrm{D0}+A^{*}(a-\mathrm{D0}), \Vert A\Vert \le C \sigma\),
and both helper norms are established. Equation (6.1) leaves the
explicit nonnegative reserve \(\Vert A p_{s} v\Vert ^{2}\). It does not
reserve another copy of the full Gauss square.

The center component of the shifted square is compared to the original
exact normal-covariant center kinetic by ordinary Young, costing
\(\epsilon  T_{\mathrm{cent}}\) plus \(\epsilon ^{-1}\) times (6.2).
Internal components are treated by the complete averaged charges in
Section \(7\); one must not apply this ordinary Young inequality to the
internal kinetic square alone.

Restore \(s\). Drop its positive norm square and expand its two cross
terms. Put \(s_{0}=-F^{-1}S_{f}\), so
\(\Vert s_{0}\Vert ^{2}\le CW_{2}\). The cross \(\langle s,a\rangle\)
differs from \(\langle s_{0},\mathrm{D0}\rangle\) by a quantity bounded
by \(\epsilon  H_{\mathrm{exc}}+C_{\epsilon} W_{2}\). To verify this,
expand \(T^{*}WT-I\) once: every factor moved onto \(\mathrm{D0}\) is
\(Z\mathrm{D0}, Z^{*}\mathrm{D0}\), or \(KK^{*}\mathrm{D0}\), all
controlled by (2.1), (2.2), (5.1). The remaining factors are uniformly
bounded. The \(n\) terms use (2.2). The leading
\(\langle s_{0},\mathrm{D0}\rangle\) is a finite \(\sum  p_{e}\) times a
Hermitian fast/slow Clifford matrix of norm \(C \kappa  /r_{e}\). Its
real form pairs with \(a_{e}\), since the Gaussian drift has purely
imaginary expectation against that Hermitian matrix. Hence it is bounded
by \(\epsilon  h_{e}+C_{\epsilon} \kappa ^{2}r_{e}^{-2}\).

For the internal part of \(\langle s,A p_{s}\rangle\), define
\(w_{a}= \sum _{e \text{ incident } a}r_{e}^{-3}\). Factoring
\(M^{-1}=R^{-1}(MR^{-1})^{-1}\) gives, for its coefficient
\(\beta _{s,a}\),

% DISPLAY h.19
% SOURCE ||β_s,av||²/w_a
% SOURCE  ≤CW2 Σ_(e incident a)r_e||Y_ev||²
% SOURCE  ≤C r_*^−3H_exc[v]+CW2[v].                              (6.3)
\begin{gather*}
\begin{aligned}
& \Vert  \beta _{s,a} v\Vert ^{2}/w_{a} \\
& \le CW_{2} \sum _{e \text{ incident } a}r_{e}\Vert Y_{e} v\Vert ^{2} \\
& \le C r_{*}^{-3}H_{\mathrm{exc}}[v]+CW_{2}[v].
\end{aligned}\tag{6.3}
\end{gather*}

Weighted Young loses \(\gamma  \sum _{a} w_{a} t_{a}\), plus
\(\gamma ^{-1}\) times the right side. This is a vanishing loss of the
complete colored internal forms because

% DISPLAY h.20
% SOURCE t_a≤H_a+c_a alpha_a,
% SOURCE Σ_a w_a alpha_a≤CκW2.                                   (6.4)
\begin{gather*}
\begin{aligned}
& t_{a}\le H_{a}+c_{a} \alpha _{a}, \\
& \sum _{a} w_{a} \alpha _{a}\le C \kappa  W_{2}.
\end{aligned}\tag{6.4}
\end{gather*}

The weights depend only on centers, so no internal derivative of a
weight occurs. Center spin cross terms use ordinary center Young. The
direct linear-Y Yukawa terms obey
\(\epsilon  H_{\mathrm{exc}}+C_{\epsilon} W_{2}\) by (3.2) with one
\(Y\) and a bounded Hermitian Clifford matrix.

All spin estimates remain on the whole colored internal module. They
make no physical-sector assumption.

\subsection{7. Complete differential-connection curvature
payment}\label{h-complete-differential-connection-curvature-payment}

Write the internal component of \(\beta  =A^{*}a\) as
\(\beta _{j}=b_{j}(Y;Z,A) \cdot  p_{F}\). This is real and fast-bosonic,
hence commutes with internal coordinates and internal Clifford matrices.
Its symmetric version is

% DISPLAY h.21
% SOURCE βhat_j=β_j−(i/2)div_F b_j.
\begin{gather*}
\begin{aligned}
& \widehat{\beta}_{j}= \beta _{j}-(i/2)\operatorname{div} _{F} b_{j}.
\end{aligned}\notag
\end{gather*}

The averaged sixteen internal charges can therefore be used with
\(p_{A}+ \widehat{\beta}\). There is no additional internal-potential
anticommutator; the only algebraic correction is a finite Clifford
contraction of the fast-bosonic curvature

% DISPLAY h.22
% SOURCE F_jk=i[P_j,P_k]=p_sym(∂_j b_k−∂_k b_j+[b_j,b_k]),
% SOURCE P_j=p_Aj+βhat_j.                                         (7.1)
\begin{gather*}
\begin{aligned}
& F_{jk}=i[P_{j},P_{k}]=p_{\mathrm{sym}}( \partial  _{j} b_{k}- \partial  _{k} b_{j}+[b_{j},b_{k}]), \\
& P_{j}=p_{Aj}+ \widehat{\beta}_{j}.
\end{aligned}\tag{7.1}
\end{gather*}

The nonsymmetric-to-symmetric kinetic conversion also has the exact
scalar correction

% DISPLAY h.23
% SOURCE (1/2)X_j(div_F b_j)+(1/4)(div_F b_j)²,
% SOURCE X_j=∂_Aj+b_j·∂Y.                                        (7.2)
\begin{gather*}
\begin{aligned}
& (1/2)X_{j}(\operatorname{div} _{F} b_{j})+(1/4)(\operatorname{div} _{F} b_{j})^{2}, \\
& X_{j}= \partial  _{Aj}+b_{j} \cdot  \partial  Y.
\end{aligned}\tag{7.2}
\end{gather*}

The companion proof Appendix \(G\) establishes the actual rootwise bound

% DISPLAY h.24
% SOURCE |<v,F_jkv>|+|<(7.2)v,v>|
% SOURCE  ≤C(σ²+r_*^−3)H_exc[v]+C W2[v].                          (7.3)
\begin{gather*}
\begin{aligned}
& \vert \langle v,F_{jk} v\rangle \vert +\vert \langle (7.2)v,v\rangle \vert \\
& \le C( \sigma ^{2}+r_{*}^{-3})H_{\mathrm{exc}}[v]+C W_{2}[v].
\end{aligned}\tag{7.3}
\end{gather*}

The proof does not infer (7.3) from a norm bound alone. Internal
differentiation of the explicit rootwise source decomposition preserves
the same-edge and triangle source classes and adds one factor at most
\(C/r_{*}\). Thus
\(\Vert r_{*} \partial  A \widehat{\beta} v\Vert ^{2}\le C[( \sigma ^{2}+r_{*}^{-3})H_{\mathrm{exc}}[v]+W_{2}[v]]\).
Weighted Cauchy with \(\Vert r_{*}^{-1}v\Vert ^{2}\le W_{2}[v]\) pays
the derivative curvature. Symmetry bounds each commutator-curvature
expectation by
\(2\Vert  \widehat{\beta}_{j} v\Vert \Vert  \widehat{\beta}_{k} v\Vert\).
The scalar divergence correction is bounded by \(CW_{2}\) from the
actual first and second coefficient derivatives. Averaged-charge
Cauchy--Schwarz therefore yields

% DISPLAY h.25
% SOURCE H_I(p+β)≥(1−ε)H_I−C_ε(σ²+r_*^−3)H_exc−C_εW2,           (7.4)
\begin{gather*}
\begin{aligned}
& H_{I}(p+ \beta  )\ge (1- \epsilon  )H_{I}-C_{\epsilon} ( \sigma ^{2}+r_{*}^{-3})H_{\mathrm{exc}}-C_{\epsilon} W_{2},
\end{aligned}\tag{7.4}
\end{gather*}

retaining the entire interacting \(H_{I}\), not a fixed fraction of
\(t_{A}\) alone.

\subsection{8. Density, vacuum defect, and the final
allocation}\label{h-density-vacuum-defect-and-the-final-allocation}

The density term is bounded by \(CW_{2}\) uniformly on the pairwise
tube. Indeed \(M R^{-1}\) and its inverse are uniformly bounded; each
slow or fast derivative of their explicit normalized entries is
\(\le C/r_{*}\), and a second derivative is \(\le C/r_{*}^{2}\). The
principal co-metric and the required normal-covariant coefficient
derivatives obey the same bounds. In
\(d^{-1}\operatorname{div} (a \partial  d)\), these give
\(C/r_{*}^{2}\le CW_{2}\). This argument uses actual normalized matrices
and does not compactify a set of gap ratios.

The frozen \(E_{0,\mathrm{total}}\) is paid separately by the incident
commutator-square result Appendix \(B\). Directly,
\(\vert E_{0,\mathrm{total}}\vert \le C \sum _{a} w_{a}V_{a}\le Cr_{*}^{-3}H_{I}+C \kappa  W_{2}\),
with fixed endpoint multiplicities absorbed in \(C\). This uses
\(V_{a}\le H_{a}+c_{a} \alpha _{a}\) and the incident bound
\(w_{a} \alpha _{a}\le C \kappa  W_{2}\). The first term is a vanishing
complete-internal-form loss; the second remains in the allowed
inverse-square remainder. Thus this local reserve does not require a
separate center-Hardy argument. \(E_{0,\mathrm{total}}\) is never
included in \(H_{\mathrm{exc}}\) or a virtual excitation denominator.

Equations (7.3), (3.1), (5.5), (6.1)--(6.4), the density bound, and the
exact Hamiltonian potential/Yukawa inventory give, for every desired
\(\delta  >0\), constants
\(\kappa _{\delta} , \sigma _{\delta} , R_{\delta}\) and \(C_{\delta}\),
independent of all gap ratios,

% DISPLAY h.26
% SOURCE H_full[v]≥(1−δ)(H_I[v]+T_cent[v]+H_exc[v])−C_δW2[v],     (8.1)
\begin{gather*}
\begin{aligned}
& H_{\mathrm{full}}[v]\ge (1- \delta  )(H_{I}[v]+T_{\mathrm{cent}}[v]+H_{\mathrm{exc}}[v])-C_{\delta} W_{2}[v],
\end{aligned}\tag{8.1}
\end{gather*}

on the tube with
\(\kappa  \le  \kappa _{\delta} , \sigma  \le  \sigma _{\delta}\) and
\(r_{*}\ge R_{\delta}\). Uniformity for nested smaller fast tubes
follows by proving the estimate at the fixed positive width
\(\sigma _{\delta}\) and restricting its core; \(C_{\delta}\) is not
obtained by substituting a varying \(\sigma\) into \(\sigma ^{-1}\). A
nonnegative fraction of \(\Vert A p_{s} v\Vert ^{2}\) and \(V_{4}\) may
also be retained by allocating their positive terms explicitly. This is
not a reserve of the original full Gauss square.

For \(v\) fiberwise orthogonal to the adapted product vacuum,
\(H_{\mathrm{exc}}[v]\ge c r_{*}\Vert v\Vert ^{2}\). Since
\(W_{2}\le (\#\text{edges})r_{*}^{-2}\), increasing \(R_{\delta}\)
absorbs \(C_{\delta} W_{2}\) into \(\delta  H_{\mathrm{exc}}\). The
resulting complementary reserve is near-unit in the full colored
\(H_{I}\), center kinetic, and shifted excitation form. Starting (8.1)
with \(\delta  /2\) gives a final displayed coefficient \(1- \delta\)
after this absorption.

For a projection-stable tube realization first fix a positive cutoff
width \(\sigma _{0}\le  \sigma _{\delta}\), then choose the exterior
radius so \(\sigma _{0}^{2}r_{*}^{3}\) is large. Choose a smooth
\(\chi _{0}\) with support strictly inside the validity tube and equal
to one on a smaller tube, using the smooth harmonic minimum of pair
separations. Put \(U_{\sigma} = \chi _{0}U/\Vert  \chi _{0}U\Vert\). The
fixed-rank Gaussian tail gives
\(\Vert U_{\sigma} -U\Vert \le  \tau  \le C \exp (-c \sigma _{0}^{2}r_{*}^{3})\),
uniformly in all gap ratios. For \(v\) fiberwise orthogonal to
\(U_{\sigma} , \vert \langle U,v\rangle \vert \le  \tau  \Vert v\Vert\),
hence \(H_{\mathrm{exc}}[v]\ge cr_{*}(1- \tau ^{2})\Vert v\Vert ^{2}\).
Both \(U_{\sigma} f\) and \(v=u-U_{\sigma} f\) have Dirichlet form
support. Thus the same complementary reserve holds for this cutoff
projection once the exterior radius is enlarged. Gaussian derivative
tails needed in a subsequent mixed-source calculation are separate
estimates.

The reserve retains the actual normal-bundle center derivative, with no
replacement by a bounded Berry coefficient. On the low vacuum line the
scalar pure-internal and mixed center/internal Berry comparison is the
rootwise full-charge lemma in source Section \(4.2\) in Appendix \(F\):
its costs are \(C_{N}( \kappa  + \epsilon ^{-1} \kappa ^{4})W_{2}\).
This does not require a flat determinant line.

No claim about mixed low/complement source estimates follows merely from
(8.1), especially if those estimates spend the full original Gauss
square instead of the actual reserve left by (6.1). The full
inverse-square Feshbach conclusion remains separate from this
complementary theorem.
